\scanpage{302}
\mspage{452}

Let us note that (33), can be written $\rho=-(1+N_1)\sigma$ i.e.
\deleted{\illegible{1}}
\[
\rho+\sigma=-N_1\sigma \tag{35}
\]
multiplying by $N_1$ on the left we get
\[
N_1\rho+N_1\sigma=0
\]
(since $N_1^2=0$), moreover
\[
N_1\rho=\rho N_0,\qquad N_1\sigma=\sigma N_0
\]
(for $N_1=\rho N_0\rho^{-1}=\sigma N_0\sigma^{-1}$), so (\unsure{35}) can be written
\[
\rho+\sigma=-N_1\sigma=+N_1\rho=-\sigma N_0=+\rho N_0
\overset{\text{definition}}{=}N . \tag{\unsure{35|36}}
\]
moreover when $\rho$ is in the general form
$\rho=\begin{pmatrix}0&u\\v&w\end{pmatrix}$ we find
\[
N\deleted{\illegible{1}}=\begin{pmatrix}0&0\\0&u\end{pmatrix},
\qquad N^2=\begin{pmatrix}0&0\\0&u^2\end{pmatrix}
\]
so the normalization (34) $u=1$ also means
\[
\rho+\sigma=\begin{pmatrix}0&0\\0&1\end{pmatrix}\deleted{\unsure{by}}, \tag{37}
\]
\inserted{or again (since $u$ is invertible)
\[
N^2=N\qquad (\text{where }N=\rho+\sigma) \tag{37 bis}
\]}
or again that $\sigma$ is given in terms of $\rho$ by
\[
\sigma=\begin{pmatrix}0&-1\\-v&1-w\end{pmatrix}. \tag{38}
\]
moreover, $N_0$ \inserted{$=\rho N_1\rho^{-1}$} will be of the form
$\begin{pmatrix}0&\alpha\\0&0\end{pmatrix}$,
an immediate calculation shows that $\alpha$ is given by the formula
$\unsure{u}=\alpha v$, i.e. (\unsure{if} we normalize by $u=1$)

\margin{\uline{N.B.} We see \unsure{at once} that the conditions considered so far
\unsure{imply} that $1,\rho,\sigma$ \unsure{are linearly independent}
\unsure{in each} fiber, i.e. $T_\rho\cap T_\sigma=\unsure{\mathbb G_m}$ !
i.e. $\widetilde T_\rho\cap\widetilde T_\sigma=\{1\}$.}

\scanpage{303}

\[
N_0=\begin{pmatrix}0&1/v\\0&0\end{pmatrix},\qquad
\varepsilon_0=\begin{pmatrix}1&1/v\\0&1\end{pmatrix}. \tag{39}
\]

\margin{\uline{N.B.} We will have
\[
\boxed{\sigma^{-1}\rho=-\varepsilon_0}\deleted{\unsure{N_0}} \tag{39 bis}
\]
For $\sigma^{-1}\rho=\sigma^{-1}(\rho\sigma^{-1})\sigma
=\sigma^{-1}(-\varepsilon_1)\sigma
=-\sigma^{-1}(\varepsilon_1)=-\varepsilon_0$.}

\deleted{Let us return}
These canonical normalizations having been made, let us return to the proposition on page 436 ; since $\varepsilon_1$ has been chosen
(in as canonical a way as can be) there is no longer any drawback to
\deleted{\illegible{1}}\inserted{leaving} equation (32) (put into the apparently more intrinsic form
(11) \unsure{or} (12)) in the \deleted{\illegible{1}} initial form,
where a priori we were thinking of introducing the parameter $\gamma$ with
\[
[\alpha,\rho]=[\beta,\rho]\varepsilon_1(\gamma)\unsure{\,;}\quad -
\]
but it appears that it is $\gamma+1=\mu$ that has a particularly remarkable significance
in the theory --- so we \unsure{keep|will keep} its form (32), with the factor $\varepsilon_1(\mu-1)$.

The expression (97) for $P(\unsure{\zeta})$, thanks to (37), can be written
\[
P(\zeta)=D\qquad\text{for }\zeta=\begin{pmatrix}A&B\\C&D\end{pmatrix}, \tag{40}
\]

\scanpage{304}
\mspage{453}

Thus the proposition tells us that the scheme \inserted{$\mathcal X_{\rho,\sigma}$}
of simultaneous solutions in $(\zeta,\eta,\mu)$ of
\[
\det\zeta=\det\eta=1,\qquad \zeta=\eta\varepsilon_1(\mu-1),\qquad
\operatorname{Tr}(\zeta\rho)=\operatorname{Tr}\rho,\qquad
\operatorname{Tr}\eta\deleted{\unsure{\zeta|\rho}}\inserted{\sigma}
=\operatorname{Tr}\unsure{\sigma|\alpha} \tag{41}
\]
is isomorphic, via $(\zeta,\eta,\mu)\longrightarrow\zeta$, to that of the
$\zeta=\begin{pmatrix}A&B\\C&D\end{pmatrix}$ satisfying the conditions
\[
\underbrace{\det\zeta}_{AD-BC}=1,\qquad
\underbrace{\operatorname{Tr}\zeta\rho}_{(C+wD)+vB}
=\underbrace{\operatorname{Tr}\rho}_{w},\qquad
\text{and }D\text{ an invertible section of }\mathcal O_S \tag{42}
\]
and we recover $\mu$ in terms of $\zeta$ by
\[
\mu=1/D \tag{43}
\]
(here $c=\operatorname{Tr}(\rho+\sigma)=\operatorname{Tr}\begin{pmatrix}0&0\\0&1\end{pmatrix}=1$),
whence $\eta$ by $\eta=\zeta\varepsilon_1(1-\mu)$.

The subgroup scheme $B'_0$ of page 442$'$ is none other than the scheme of matrices
\[
B'_0=\left\{\begin{pmatrix}\mu&\lambda\\0&1\end{pmatrix}
\,\middle|\,\mu\in\Gamma\mathbb G_m,\ \lambda\in\Gamma\mathbb G_a\right\} \tag{44}
\]
We are going to make explicit the isomorphism (79)
\[
\underline u\longmapsto\bigl([\underline u,\rho],[\underline u,\sigma],\det\underline u\bigr) \tag{45}
\]
in terms of the coefficients $\mu,\lambda$ of $B'_0$, with the help of the parameters
$v,w$ in $\rho=\begin{pmatrix}0&1\\v&w\end{pmatrix}$.

\scanpage{305}

We find
\deleted{\unsure{$\underline u(\rho)=$}}
\[
\begin{aligned}
{}[\underline u,\rho]
&=\begin{pmatrix}\mu&\lambda\\0&1\end{pmatrix}
\begin{pmatrix}0&1\\v&w\end{pmatrix}
\frac1\mu\begin{pmatrix}1&-\lambda\\0&\mu\end{pmatrix}
\left(-\frac1v\right)\begin{pmatrix}w&-1\\-v&0\end{pmatrix}\\
&=-\frac1{\mu v}\begin{pmatrix}\lambda v&\mu+\lambda w\\v&w\end{pmatrix}
\begin{pmatrix}w+\lambda v&-1\\-\mu v&0\end{pmatrix}\\
&=-\frac1{\mu v}
\begin{pmatrix}
\deleted{\illegible{2}\unsure{-\mu)vw+\lambda}}\inserted{v^2\lambda^2-v\mu^2-vw\lambda\mu+vw\lambda}&-\lambda v\\
v^2\lambda-vw\mu+vw&-v
\end{pmatrix}\\
&=\frac1\mu\begin{pmatrix}
-v\lambda^2+w\lambda\mu+\mu^2-w\lambda&\unsure{\lambda}\\
-v\lambda+w\mu-w&1
\end{pmatrix}
\end{aligned}
\]

\margin{\unsure{now|and} for $\sigma=\begin{pmatrix}0&-1\\v'&w'\end{pmatrix}$
with $v'=-v$, $w'=1-w$}

\[
\begin{aligned}
{}[\underline u,\sigma]
&=\begin{pmatrix}\mu&\lambda\\0&1\end{pmatrix}
\begin{pmatrix}0&-1\\v'&w'\end{pmatrix}
\frac1\mu\begin{pmatrix}1&-\lambda\\0&\mu\end{pmatrix}
\left(\frac1{v'}\right)\begin{pmatrix}w'&1\\-v'&0\end{pmatrix}\\
&=\frac1{\mu v'}\begin{pmatrix}\lambda v'&-\mu+\lambda w'\\v'&w'\end{pmatrix}
\begin{pmatrix}w'+\lambda v'&1\\-\mu w'&0\end{pmatrix}\\
&=\frac1{\mu v'}\begin{pmatrix}
v'^2\lambda^2+v'\mu^2-v'w'\lambda\mu+v'w'\lambda&\lambda v'\\
v'^2\lambda-v'w'\mu+v'w'&v'
\end{pmatrix}\\
&=\frac1\mu\begin{pmatrix}
v'\lambda^2\unsure{-}\deleted{\unsure{\mu}}w'\lambda\mu+\mu^2+w'\lambda&\lambda\\
v'\lambda-w'\mu+w'&1
\end{pmatrix}\\
&=\frac1\mu\begin{pmatrix}
-v\lambda^2+(w-1)\lambda\mu+\mu^2\unsure{-}(w-1)\lambda&\lambda\\
-v\lambda+(w-1)\mu-(w-1)\deleted{\unsure{\lambda}}&1
\end{pmatrix}
\end{aligned}
\]

In summary
\[
\left\{\begin{aligned}
\zeta=[\underline u,\rho]
&=\frac1\mu\begin{pmatrix}
\mu^2+(-v\lambda+w(\mu\unsure{-}1))\lambda&\lambda\\
-\deleted{\unsure{(}}v\lambda+w(\mu-1)&1
\end{pmatrix}\\
\eta=[\underline u,\sigma]
&=\frac1\mu\begin{pmatrix}
\mu^2+(-v\lambda+(w-1)(\mu-1))\lambda&\lambda\\
-v\lambda+(w-1)(\mu-1)&1
\end{pmatrix}\quad\text{and, of course}\\
\mu&=\det\underline u\ \deleted{\unsure{=\mu}}
\end{aligned}\right. \tag{46}
\]
so we have, if $\zeta=\begin{pmatrix}A&B\\C&D\end{pmatrix}$
\[
B=\lambda/\mu,\quad D=1/\mu\qquad\text{i.e.} \tag{47}
\]
\[
\mu=1/D,\quad\lambda=B/D,\qquad\text{\unsure{i.e.|of which} }\underline u=\begin{pmatrix}1/D&B/D\\0&1\end{pmatrix} \tag{48}
\]

\scanpage{306}
\mspage{453bis}

Likewise, if $\eta=\begin{pmatrix}R&S\\T&U\end{pmatrix}$, we will have \deleted{$S=B,U=D$}
\deleted{whence}
\[
\left\{\begin{aligned}
S=B&=\lambda/\mu,\quad U=D=1/\mu\qquad\text{whence}\\
\deleted{\unsure{and}}\ \mu&=1/U,\quad\lambda=S/U
\end{aligned}\right. \tag{47 bis}
\]
so
\[
\underline u=\begin{pmatrix}1/U\deleted{\unsure{(=\mu)}}&S/U\\0&1\end{pmatrix}. \tag{\unsure{48 bis}}
\]

Let us now consider the equation
\[
[\alpha,\rho]=[\beta,\sigma]\varepsilon_1(\mu-1),\qquad\text{with} \tag{49}
\]
\[
\alpha=\begin{pmatrix}a&b\\c&d\end{pmatrix},\quad
\beta=\begin{pmatrix}r&s\\t&u\end{pmatrix}, \tag{\unsure{50}}
\]
and set, for more \unsure{elegant} calculations
\[
\rho=\begin{pmatrix}0&\unsure{v_1}\\\unsure{v}&w\end{pmatrix},\quad
\sigma=\begin{pmatrix}0&v'_1\\v'&w'\end{pmatrix} \tag{51}
\]
(we have written $v_1,v'_1$ instead of $u,u'$, so as not to \unsure{interfere} with the notation $u$ for one of the coefficients of $\beta$), with
\[
v'=-v,\quad v'_1=-\unsure{v_1|v'},\quad w+w'=v_1,\qquad\text{let us set} \tag{52}
\]
\[
\delta_\alpha=\det\alpha=ad-bc,\quad
\delta_\beta=\det\beta=ru-st,
\tag{53}
\]

\margin{We do not set $v_1=1$, so as \unsure{not to have} to \unsure{redo} the same calculation for $[\beta,\sigma]$ as for $[\alpha,\rho]$ !}

we then find
\[
\begin{aligned}
\zeta=[\alpha,\rho]
&=\begin{pmatrix}a&b\\c&d\end{pmatrix}
\begin{pmatrix}0&v_1\\v&w\end{pmatrix}
\frac1{\delta_\alpha}\begin{pmatrix}d&-b\\-c&a\end{pmatrix}
\deleted{\unsure{(}}\frac1{-vv_1}\begin{pmatrix}w&-v_1\\-v&0\end{pmatrix}\\
&=\unsure{-}\frac1{\delta_\alpha vv_1}
\begin{pmatrix}bv&av_1+bw\\dv&cv_1+dw\end{pmatrix}
\begin{pmatrix}dw+bv&-dv_1\\-cw-av&cv_1\end{pmatrix}
=\begin{pmatrix}A&B\\C&D\end{pmatrix}
\end{aligned}
\]
with
\[
\left\{\begin{aligned}
(-\delta_\alpha vv_1)A&=\ldots\quad\text{(a bore).},
&(-\delta_\alpha vv_1)B&=-vv_1bd+v_1^2ac+v_1wbc\\
(-\delta_\alpha vv_1)C&=\ldots\quad\text{(a bore)},
&(-\delta_\alpha vv_1)D&=-vv_1d^2+v_1wcd+v_1^2c^2
\end{aligned}\right. \tag{54}
\]
\margin{N.B. the formulas for $B,D$ \unsure{simplify} by $v_1$}

\scanpage{307}

and we likewise find
\[
\eta=\unsure{(\beta,\sigma)|[\beta,\sigma]}=
\begin{pmatrix}R&S\\T&U\end{pmatrix}
\]
with
\[
\tag{55}
\begin{array}{ll}
(-vv_1\delta_\beta)R=\ldots
&(-vv_1\delta_\beta)S
 =(v_1v)su+v_1^{2}tr-v_1w'st\\
(-vv_1\delta_\beta)T=\ldots
&(-vv_1\delta_\beta)U
 =-(v_1v)u^2-v_1w'tu+v_1^{2}t^2.
\end{array}
\]
\margin{\uline{N.B.} we can divide by $v_1$.}

\unsure{We} deduce from this in particular
\deleted{\illegible{1}}
\[
\tag{56}
\left\{
\begin{aligned}
D&=\frac{-v\deleted{\unsure{v_1}}d^2+
\deleted{\unsure{v_1}}wcd+v_1^{\deleted{2}}c^2}
{-v\deleted{\unsure{v_1}}(ad-bc)}\\
\mu=1/D&=\frac{-v\deleted{\unsure{v_1}}(ad-bc)}
{-v\deleted{\unsure{v_1}}d^2+
\deleted{\unsure{v_1}}wcd+v_1^{\deleted{2}}c^2}\\
\deleted{(57)}\quad\lambda&=\deleted{\unsure{C/D}}\inserted{B/D}
=\frac{-v\deleted{\unsure{v_1}}bd+v_1^{\deleted{2}}ac+
\deleted{\unsure{v_1}}wbc}
{-v\deleted{\unsure{v_1}}d^2+
\deleted{\unsure{v_1}}wcd+v_1^{\deleted{2}}c^2}.
\end{aligned}
\right.
\]
which thus makes it possible to \unsure{make explicit}
\[
\underline{u}=\begin{pmatrix}1/D=\mu&B/D=\lambda\\0&1\end{pmatrix}
\]
in terms of $a,b,c,d$ (and the parameters $u_1,v,w$
which determine $\rho$).

For \uline{$v_1=1$}, the formula for $D$ in (56), can be written
\[
\tag{57}
D=\det(d\rho+c)/-v\delta_\alpha
=\det(d\rho+c)/\det\rho\det\alpha
\]
\deleted{the condition \unsure{to} \illegible{1} on $\alpha$ ``$D$ invertible''}
\[
\tag{58}
D(\alpha)\text{ invertible}\quad\Longleftrightarrow\quad
d\rho+c\text{ invertible}.
\]
but I realize that this is a \unsure{flat} form,